%! TeX program = lualatex
\documentclass[../main.tex]{subfiles}

\usetikzlibrary{intersections,calc}

\newenvironment{solution}{
  \begin{proof}[Sample solution]
    \color{main}
  }{
  \end{proof}
}

\begin{document}

\chapter*{Week 12 practice problems}

\setcounter{chapter}{12}

\begin{exercise}
  Suppose \(A\) is a \(2\)-by-\(2\) matrix and \(\vec{v} = \begin{bmatrix} 2 \\ 1 \end{bmatrix}\) is an eigenvector for \(A\).  Sketch \(\vec{v}\) and \(A\vec{v}\) for each of the scenarios below.
  \begin{enumerate}
    \item The eigenvalue for \(\vec{v}\) is larger than \(1\).
    \item The eigenvalue for \(\vec{v}\) is exactly \(1\).
    \item The eigenvalue for \(\vec{v}\) is between \(0\) and \(1\), exclusive.
    \item The eigenvalue for \(\vec{v}\) is negative.
  \end{enumerate}

  % \begin{tikzpicture}[scale=1]
  %   \begin{axis}[basic curve, xmin=-5, xmax=5,ymin=-5, ymax=5, width=5in, xlabel=\(x_{1}\), ylabel=\(x_{2}\)]
  %     \addplot coordinates {(0,0)};
  %   \end{axis}
  % \end{tikzpicture}
\end{exercise}

\begin{exercise}
  Suppose \(\vec{x}(t+1) = A \vec{x}(t)\) models some population dynamics. Suppose further that \(\vec{u}\) and \(\vec{v}\) are two eigenvectors for \(A\). 
  \begin{enumerate}
    \item If \(A \vec{u}\) is longer than \(\vec{u}\), and \(A \vec{v}\) is shorter than \(\vec{v}\), does the population grow unbounded, stabilize, or go extinct in the long run?
    \item If \(A \vec{u}\) is a little bit shorter than \(\vec{u}\), and \(A \vec{v}\) is a lot shorter than \(\vec{v}\), does the population grow unbounded, stabilize, or go extinct in the long run?
  \end{enumerate}
\end{exercise}

\begin{exercise}
  Let \(x_{1}(t)\) be the number of juveniles in year \(t\), and \(x_{2}(t)\) be the number of adults in year \(t\).  The per-capita birth rate among adults is \(1/4\). The survival rate among adults is \(1/4\). Lastly, \(1/2\) of the juveniles grows into adults in the next year. 

  \begin{enumerate}
    \item Sketch the transition diagram of this population.
    \item Model the population dynamics using a (vector) recurrence.
    \item Solve the recurrence, assuming that \(x_{1}(0) = 32\) and \(x_{2}(0) = 16\).
    \item Determine the long term behaviour of the population. Does the population grow unbounded, stabilize or go extinct? If they don't go extinct, what is the ratio of \(x_{2}(t) / x_{1}(t)\) as \(t\) approaches infinity?
  \end{enumerate}
\end{exercise}

\begin{exercise}
  A traffic light at a quiet intersection is broken.  At any given time, only the green and the red lights up.
  \begin{itemize}
    \item If it's red at time \(t\), then there is a \(20\%\) chance it will turn green at time \(t+1\).
    \item If it's green at time \(t\), then there is a \(50\%\) chance it will turn red at time \(t+1\).
  \end{itemize}

  \begin{enumerate}
    \item Sketch the transition diagram of the traffic light states.
    \item Model the state transitions using a (vector) recurrence.
    \item Solve the recurrence, assuming that \(x_{1}(0) = 0.5\) and \(x_{2}(0) = 0.5\).
    \item Determine \(x_{1}(t)\) and \(x_{2}(t)\) as \(t\) approaches infinity.
  \end{enumerate}
\end{exercise}

\begin{exercise}
  Predict the long-term behaviour of the model \(\vec{x}(t+1) = \begin{bmatrix} 2 & -3/2 \\ -12 & -1\end{bmatrix} \vec{x}(t)\) with \(x_{1}(0) = 3\) and \(x_{2}(0) = 6\). 
\end{exercise}

\begin{exercise}
  Consider the recurrence \(\vec{x}(t+1) = \begin{bmatrix} 5 & 4 \\ 1 & 2\end{bmatrix} \vec{x}(t)\). If \(\mu = 1\) and \(\lambda = 6\) are its eigenvalues, find the ratio \(x_{2}(t) / x_{1}(t)\) as \(t\) approaches infinity.
\end{exercise}

\begin{exercise}
  Solve the initial-value problem \( \frac{d\vec{x}}{dt} = \begin{bmatrix} 3 & 0  \\ 0 & -2 \end{bmatrix} \vec{x}(t) \) with \(\vec{x}(0) = \begin{bmatrix} 1/2 \\ 1/3 \end{bmatrix}\).
\end{exercise}

\begin{exercise}
  Find the general solutions for
  \(
  \begin{array}{rcrcr}
    x'_{1}(t) &=& -2 x_{1}(t) &+&  x_{2}(t) \\
    x'_{2}(t) &=&    x_{1}(t) &+& 2x_{2}(t)
  \end{array}
  \).
\end{exercise}

\begin{exercise}
  Find the specific solutions for
  \(
  \begin{array}{rcrcr}
    u'_{1}(t) &=& 2 u_{1}(t) &+& 2u_{2}(t) \\
    u'_{2}(t) &=&   u_{1}(t) &+& 3u_{2}(t)
  \end{array}
  \) with \(
  \begin{array}{rcr}
    u_{1}(0) &=&  1 \\
    u_{2}(0) &=&  3 \\
  \end{array}
  \).
\end{exercise}
\end{document}
